\begin{frame}[t]{跨模态继续预训练：将蛋白丰度接入 RNA 模型}
\tagline{Burq et al.，bioRxiv v1，2026-08-19；模型规模 70M（约 7,000 万参数）。}
\centering
\begin{tikzpicture}[n/.style={box,text width=3.55cm,minimum height=1.18cm,inner sep=8pt}]
\node[n] (a) at (0,0) {\textbf{已有 Tahoe-x1 70M}\\\scriptsize 加载 RNA 预训练权重};
\node[n,fill=PaleTeal] (b) at (4.7,0) {\textbf{蛋白组继续预训练}\\\scriptsize 48,843 样本，1 epoch};
\node[n] (c) at (9.4,0) {\textbf{通用分子表示}\\\scriptsize 基因级 / 样本级下游评估};
\draw[flow] (a)--(b);\draw[flow] (b)--(c);
\end{tikzpicture}
\vspace{.20cm}
\begin{columns}[T,onlytextwidth]
\begin{column}{.48\textwidth}
\subhead{统一身份与数值编码}
\bodyfont 蛋白 $\rightarrow$ 对应基因 token；丰度 $\rightarrow$ expression value。每个样本内部做分位数分箱，减轻不同来源强度尺度的差异。
\end{column}
\begin{column}{.48\textwidth}
\subhead{沿用自监督训练方案}
\bodyfont 采用原 Tahoe-x1 training recipe；展示原始与批次校正两种视图，称为 format-aware augmentation。评估用药物扰动样本从继续预训练中排除。
\end{column}
\end{columns}
\remark{\textbf{不要求 RNA--蛋白样本配对，也未按相关性聚类。}连接来自共同基因身份与兼容输入格式。}
\source{[25] §2.2。沿用已有 checkpoint，而非本研究重新从头训练 RNA 模型；70M = 约 7,000 万参数。论文未完整列出本轮优化器、学习率和 masking 配置。}
\end{frame}
